%%%%%%%%%%  scan idx 592  |  folio 579  %%%%%%%%%%
\SGAPrintPage{579}%
\(K^\bullet(T)\) and the degree filtration on the second member follows readily from (V 6.6.4) and from the fact that the \(c^i\) define a \emph{Chern theory} of \(K^\bullet(T)\) in the sense of (V 6.13), when \(K^\bullet(T)\) is regarded as a \(\lambda\)-algebra augmented over the binomial ring \(H^0(T,\mathbb Z)\)). One thereby obtains maps
\sgaStableTarget{sga6.e10.eq.7-13-7}\begin{equation}\tag{7.13.7}
  \operatorname{c\ell}:\Gr^i(T)\longrightarrow H^{2i}(T,\mathbb Q_\ell(i))
  \stackrel{\mathrm{dfn}}{=}
  H^{2i}(T,\mathbb Z_\ell(i))\otimes_{\mathbb Z_\ell}\mathbb Q_\ell,
\end{equation}
by the formula
\sgaStableTarget{sga6.e10.eq.7-13-8}\begin{equation}\tag{7.13.8}
  \operatorname{c\ell}(x^{(i)})=
  (1/(-1)^{i-1}(i-1)!)\,\psi_i(x^{(i)}),
\end{equation}
compare (XIV 5.1 and 5.2). With this understood, the commutativity of diagrams (7.13.2) and (7.13.4) trivially implies that of the analogous diagrams
\sgaStableTarget{sga6.e10.eq.7-13-9}\begin{equation}\tag{7.13.9}
\begin{tikzcd}[column sep=large,row sep=large]
\Gr^i(X_t)_{\mathbb Q} \arrow[r,"\operatorname{c\ell}"] \arrow[d,"\sigma_{\mathbb Q}"'] &
H^{2i}(X_t,\mathbb Q_\ell(i))\\
\Gr^i(X_s)_{\mathbb Q} \arrow[r,"\operatorname{c\ell}"] &
H^{2i}(X_s,\mathbb Q_\ell(i)) \arrow[u,"\sigma^{-1}"]
\end{tikzcd}
\;,
\end{equation}
\sgaStableTarget{sga6.e10.eq.7-13-10}\begin{equation}\tag{7.13.10}
\begin{tikzcd}[column sep=large,row sep=large]
\Gr^i(X_{\bar t})_{\mathbb Q} \arrow[r,"\operatorname{c\ell}"] \arrow[d] &
H^{2i}(X_{\bar t},\mathbb Q_\ell(i)) \arrow[d,"\wr"]\\
\Gr^i(X_{\bar s})_{\mathbb Q} \arrow[r,"\operatorname{c\ell}"] &
H^{2i}(X_{\bar s},\mathbb Q_\ell(i))
\end{tikzcd}
\;.
\end{equation}
\sgaStableTarget{sga6.e10.rem.7-13-11}\textbf{Remark 7.13.11.} The commutativity of the last two diagrams, which can be expressed by saying that specialization of algebraic cycles is compatible with formation of the associated \(\ell\)-adic cohomology classes, is not very satisfactory in the form stated, since we have had to replace the coefficients \(\mathbb Z_\ell(i)\) by \(\mathbb Q_\ell(i)\), that is, to neglect torsion. To obtain a statement that does not neglect torsion, it seems that one

%%%%%%%%%%  scan idx 593  |  folio 580  %%%%%%%%%%
\SGAPrintPage{580}%
will be forced to work with the Chow ring, since it is on that ring that the homomorphism \(\operatorname{c\ell}\) is naturally defined as a homomorphism with values in \(H^{2*}(-,\mathbb Z_\ell(*))\); compare (XIV 4.1). Unfortunately, no theory of specialization for the Chow ring has yet been developed (cf. 7.14.) that would allow one to write commutative squares analogous to (7.13.9) and (7.13.10) in terms of Chow rings and cohomology with coefficients \(\mathbb Z_\ell(*)\). All one can say when \(f:X\to S\) is smooth and proper is that there is a commutative square
\sgaStableTarget{sga6.e10.eq.7-13-12}\begin{equation}\tag{7.13.12}
\begin{tikzcd}[column sep=large,row sep=large]
Z^i(X_t) \arrow[r,"\operatorname{c\ell}"] \arrow[d,"\sigma"'] &
H^{2i}(X_t,\mathbb Z_\ell(i))\\
Z^i(X_s) \arrow[r,"\operatorname{c\ell}"] &
H^{2i}(X_s,\mathbb Z_\ell(i)) \arrow[u]
\end{tikzcd}
\;,
\end{equation}
analogous to (7.13.9), and a corresponding square analogous to (7.13.10), which we leave to the reader to write down, where \(Z^i\) denotes the group of algebraic cycles (modulo no relation at all) of codimension \(i\), where the horizontal arrows are the ``associated cohomology class'' homomorphisms defined in (SGA 5 IV), and where \(\sigma\) denotes the specialization homomorphism for algebraic cycles, associating with a cycle \(Z_t\) the cycle \(Z_s\), where \(Z\) is the schematic closure of \(Z_t\) in \(X\) (cf. 7.14). This commutativity follows from the following fact, established in (SGA 5 IV): if \(Z\) is a closed subscheme of \(X\) of codimension \(\geq i\) on every fiber, one can associate with it a cohomology class \(\operatorname{c\ell}(Z)\in H^{2i}(X,\mathbb Z_\ell(i))\) whose formation commutes with every base change and, in particular, such that its restriction to \(X_t\) (resp. \(X_s\)) equals \(\operatorname{c\ell}(Z_t)\) (resp. \(\operatorname{c\ell}(Z_s)\)). Of course, when one knows that the homomorphism \(\sigma\) in (7.13.12) carries cycles rationally equivalent to zero to cycles rationally equivalent to zero (\(X\) being projective over \(S\)),

%%%%%%%%%%  scan idx 594  |  folio 581  %%%%%%%%%%
\SGAPrintPage{581}%
one can replace the \(Z^i\) by the \(A^i\) in this diagram, and obtain a diagram of the desired kind in terms of the Chow rings of \(X_t\) and \(X_s\).

\sgaStableTarget{sga6.e10.rem.7-13-13}\textbf{Remark 7.13.13.} Let \(T\) be a smooth projective scheme over an algebraically closed field. It follows trivially from the formalism developed in (SGA 5 IV) that, if \(Z\) is a cycle of codimension \(i\) on \(T\) whose image in \(H^{2i}(T,\mathbb Q_\ell(i))\) is zero, then \(Z\) is numerically equivalent to zero. One of the standard conjectures in the theory of algebraic cycles asserts that the converse is also true. Moreover, by (VII 4.11), it is immediate that one obtains an equivalent conjecture by replacing the group of algebraic cycles with the group \(\Gr^\bullet(T)_{\mathbb Q}\), using definition (7.13.8), or again with the group \(K^\bullet(T)\), by stating that numerical equivalence to zero of an element of \(K^\bullet(T)\) is equivalent to vanishing of its augmentation and of its Chern classes in the \(H^{2i}(T,\mathbb Q_\ell(i))\).

With this understood, it follows immediately from the commutativity of (7.13.10) that, if \(X_{\bar t}\) satisfies the preceding conjecture, then the specialization homomorphism \(\Gr^\bullet(X_{\bar t})\longrightarrow \Gr^\bullet(X_{\bar s})\) (and, a fortiori, \(\Gr^\bullet(X_t)\longrightarrow \Gr^\bullet(X_s)\)) is compatible with numerical equivalence, as conjectured in 7.9. Note also that the conjecture invoked for this purpose (numerical equivalence implies \(\mathbb Q_\ell\)-cohomological equivalence) is weaker than the one considered in 7.12.5.

\sgaStableTarget{sga6.e10.sec.7.14}\subsection*{7.14. Specialization homomorphism for groups of cycles and cycle classes.}
If \(T\) is a noetherian scheme, denote by \(Z_i(T)\) (resp. \(Z^i(T)\)) the group of algebraic cycles on \(T\) purely of dimension \(i\) (resp. purely of codimension \(i\)). When \(f:X\to S\) is of finite type, one then defines

%%%%%%%%%%  scan idx 595  |  folio 582  %%%%%%%%%%
\SGAPrintPage{582}%
a specialization homomorphism:
\sgaStableTarget{sga6.e10.eq.7-14-1}\begin{equation}\tag{7.14.1}
  \sigma_i:Z_i(X_t)\longrightarrow Z_i(X_s),
\end{equation}
by defining, for every prime cycle \(Z_t\) of \(X_t\) purely of dimension \(i\), identified with an integral subscheme of \(X_t\), the element \(\sigma_i(Z_t)\) of \(Z_i(X_s)\) as the cycle associated with the closed subscheme \(Z_s\) of \(X_s\), where \(Z\) is the schematic closure of \(Z_t\) in \(X\) (which is a closed subscheme of \(X\) flat over \(S\) (EGA IV 2.8), with generic fiber equal to the original \(Z_t\), and satisfying \(\dim Z_s=\dim Z_t\) (EGA IV 7.1.13)). When \(X_s\) is locally equidimensional (i.e., all its irreducible components through a given point have the same dimension), one immediately verifies that the specialization homomorphism just defined induces homomorphisms
\sgaStableTarget{sga6.e10.eq.7-14-1bis}\begin{equation}\tag{7.14.1 bis}
  \sigma^i:Z^i(X_t)\longrightarrow Z^i(X_s).
\end{equation}

Now suppose that \(f\) is quasi-projective, so that the notion of rational equivalence of cycles is available on \(X_t\) and on \(X_s\) (cf. [5], where the theory is developed over an algebraically closed field). It seems very plausible that homomorphism (7.14.1) carries cycles rationally equivalent to zero to cycles rationally equivalent to zero, and thus induces a corresponding homomorphism on the ``Chow'' groups of cycle classes
\sgaStableTarget{sga6.e10.eq.7-14-2}\begin{equation}\tag{7.14.2}
  \sigma_i:A_i(X_t)\longrightarrow A_i(X_s),
\end{equation}
or again, when \(X_s\) is locally equidimensional:
\sgaStableTarget{sga6.e10.eq.7-14-2bis}\begin{equation}\tag{7.14.2 bis}
  \sigma^i:A^i(X_t)\longrightarrow A^i(X_s).
\end{equation}

One might try to prove this by a method paralleling that of 7.12, but in which the recourse to resolution of singularities and the reduction

%%%%%%%%%%  scan idx 596  |  folio 583  %%%%%%%%%%
\SGAPrintPage{583}%
to an algebraically closed residue field would become unnecessary, since \(C=\mathbb P^1_S\) is available here as a ``parameter curve'' along which to deform cycles. A difficulty nevertheless arises, for with the notation of 7.12, \(E\) must now denote a cycle on \(P^1_S\times_SX\) such that \(E_t(x_{1t})-E_t(x_{0t})\) ``is defined,'' but it by no means follows that the cycle \(E_s(x_{1s})-E_s(x_{0s})\) is also defined. As was to be expected, one is thus reduced to formulating a statement of the type of Chow's ``moving lemma'' in order to define (7.14.2).

Observe, however, that the existence of homomorphism (7.4.1) shows that (7.14.1) is compatible with the equivalence relation on cycles defined by the surjective homomorphism \(Z_\bullet(T)\longrightarrow \Gr_\bullet(T)\), which is only slightly coarser than linear equivalence; thus, when \(Z\) is quasi-projective and smooth over a field, we point out elsewhere that it differs from linear equivalence only up to torsion (XIV 4.1, 4.2). In this case one therefore obtains a homomorphism analogous to (7.14.2), or better to (7.14.2 bis), but with both members tensored by \(\mathbb Q\).

The most natural way to define a homomorphism (7.14.2) directly would seem to be to proceed as in 7.3. This would, however, require a theory of Chow groups \(A_i\) (or of the Chow ring in the regular case) for schemes more general than schemes quasi-projective over a field, including \(X\), and a contravariant law for these groups that would allow one to define the restriction homomorphism to \(X_s\), \(A_i(X)\longrightarrow A_i(X_s)\). Compare XIV 8. Once such a sufficiently general theory of Chow groups and rings is available, it should be possible to transcribe into Chow theory all the results enumerated above for the specialization homomorphism in \(K\)-theory.

